\section{Constructing the Circular Arc of a Cube Edge}
Let a finite cube edge have projected endpoints
\begin{equation}
  \vb q_1=(x_1,y_1),
  \qquad
  \vb q_2=(x_2,y_2),
\end{equation}
and let $\vb v=(x_3,y_3)$ be the inside vanishing point of its direction family.  The supporting circle of the edge is taken to be the circumcircle through these three points.

Define
\begin{equation}
  D=2\left[
  x_1(y_2-y_3)+x_2(y_3-y_1)+x_3(y_1-y_2)
  \right].
  \label{eq:circumcircle-determinant}
\end{equation}
If $|D|$ is below a tolerance, the three points are effectively collinear and the edge is drawn as a straight segment.  Otherwise the circle center is
\begin{align}
  c_x={}&\frac{
  (x_1^2+y_1^2)(y_2-y_3)
  +(x_2^2+y_2^2)(y_3-y_1)
  +(x_3^2+y_3^2)(y_1-y_2)}{D},\\
  c_y={}&\frac{
  (x_1^2+y_1^2)(x_3-x_2)
  +(x_2^2+y_2^2)(x_1-x_3)
  +(x_3^2+y_3^2)(x_2-x_1)}{D},
  \label{eq:circumcenter}
\end{align}
and its radius is
\begin{equation}
  r_c=\sqrt{(x_1-c_x)^2+(y_1-c_y)^2}.
  \label{eq:circumradius}
\end{equation}
This construction guarantees that the full supporting circle contains both finite endpoints and the relevant vanishing point.

The visible cube edge is only the portion between $\vb q_1$ and $\vb q_2$.  The implementation normally draws the shorter of the two arcs joining these endpoints.  The vanishing point constrains the curvature of the supporting circle but generally lies on an extension rather than on the visible finite segment.

\section{Choosing and Drawing the Arc Branch}
The angular coordinate of a planar point $\vb q$ around the circumcenter is
\begin{equation}
  \alpha(\vb q)=
  \operatorname{atan2}(q_y-c_y,q_x-c_x).
\end{equation}
Let
\begin{equation}
  \alpha_1=\alpha(\vb q_1),
  \quad
  \alpha_2=\alpha(\vb q_2),
  \quad
  \alpha_v=\alpha(\vb v).
\end{equation}
There are two arcs between $\alpha_1$ and $\alpha_2$.  For a finite cube edge, choose the local or shorter branch.  For a guide extension that must pass through a vanishing point, choose the clockwise or counterclockwise branch whose angular interval contains $\alpha_v$.

Once the branch is fixed, sample it as
\begin{equation}
  \vb q(s)=\vb c+r_c
  \bigl(\cos\alpha(s),\sin\alpha(s)\bigr),
  \qquad 0\leq s\leq1,
  \label{eq:arc-sampling}
\end{equation}
where $\alpha(s)$ interpolates continuously along the selected interval, with the appropriate $2\pi$ wrap-around.

A useful adaptive segment count is
\begin{equation}
  N_{\rm seg}=\max\left[
  N_{\min},
  \left\lceil\frac{r_c|\Delta\alpha|}{\ell_{\max}}\right\rceil
  \right],
  \label{eq:adaptive-segments}
\end{equation}
where $\ell_{\max}$ is the desired maximum arc length per rendered segment.  A fixed upper cap prevents very large circumcircles from producing excessive geometry.

\section{Drawing Guide Arcs toward Both Vanishing Points}
For each axis family, the renderer draws pale extensions from the finite edge toward $\vb v_{\rm in}$ and $\vb v_{\rm out}$.  The inside and outside extensions are handled separately:
\begin{enumerate}
  \item order the two projected vertices by their distance from $\vb v_{\rm in}$;
  \item construct an arc through the nearer vertex, the farther vertex, and $\vb v_{\rm in}$ for the inward continuation;
  \item construct the corresponding outward continuation using $\vb v_{\rm out}$ and the reversed endpoint ordering;
  \item choose the arc branch that actually passes through the requested vanishing point;
  \item place guide arcs behind the solid cube edges in rendering depth.
\end{enumerate}

Because the outside vanishing point is generated by Eq.~\eqref{eq:outside-vp}, the two guide pieces need not be portions of one globally identical Euclidean circle.  The current implementation prioritizes a consistent visual continuation on both sides of the finite edge.  If a single exact supporting curve is required, the preferable procedure is to sample the full three-dimensional line, map every sample with $F$, and render the resulting polyline or a fitted spline.

\section{Why the Circumcircle Approximation Is Reasonable}
The circumcircle construction is motivated by two exact perspective constraints.  First, a spatial line traces one great circle of viewing directions.  Second, all parallel lines share the same pair of limiting directions and hence the same vanishing points.  Requiring a projected edge to lie on a circle through its two finite endpoints and the appropriate vanishing point imposes both constraints in a simple form.

For stereographic projection, the resulting circle is exact.  For the active circular-coordinate map $F$, however, the complete mapped great circle is not analytically proven to be a Euclidean circle.  The code therefore uses the circumcircle as a geometrically constrained interpolation.  Its accuracy can be tested directly.  Sample points along the original edge,
\begin{equation}
  \vb P(s)=(1-s)\vb P_1+s\vb P_2,
  \qquad 0\leq s\leq1,
\end{equation}
map them with $F$, and calculate the circle residual
\begin{equation}
  \delta(s)=
  \left|
  \lVert F(\vb P(s))-\vb c\rVert-r_c
  \right|.
  \label{eq:circle-residual}
\end{equation}
If $\max_s\delta(s)$ is small compared with a pixel or the rendered line width, the circular approximation is visually adequate.  Otherwise the sampled mapped edge should be rendered directly.

\section{Practical Implementation Sequence}
A complete cube renderer can be organized as follows.
\begin{enumerate}
  \item Transform the eight local cube vertices into world coordinates.
  \item For every vertex, compute the unit viewing direction from the eye.
  \item Evaluate $\theta$ and $\phi$, construct the two circular coordinate lines, intersect them, and choose the valid point inside the boundary disk.
  \item Obtain the three transformed cube-axis directions from representative edges.
  \item For each axis, choose the forward orientation, map it to $\vb v_{\rm in}$, and construct $\vb v_{\rm out}$ from Eq.~\eqref{eq:outside-vp}.
  \item Associate each of the twelve edges with its $x$, $y$, or $z$ direction family.
  \item For each edge, fit the circumcircle through its two projected endpoints and the corresponding inside vanishing point.
  \item Draw the finite edge using the shorter endpoint-to-endpoint arc.
  \item Draw guide extensions toward the inside and outside vanishing points using branch-aware arcs.
  \item Replace unstable or nearly collinear constructions by straight lines.
\end{enumerate}

In the repository, the relevant routines are in \texttt{js/projections/hemispherical-projection.js}.  The cube edge list and the mapping from edge pairs to axis families are stored in \texttt{js/config.js}, while the transformed world vertices are assembled in \texttt{js/utils/three-utils.js}.

\section{Numerical and Geometric Edge Cases}
Several special cases are important in an interactive implementation.
\begin{itemize}
  \item \emph{Central vanishing point:} if $\lVert\vb v_{\rm in}\rVert$ is below a tolerance, draw the direction family as radial straight lines.
  \item \emph{Boundary vanishing point:} if $|d_z|$ is below a tolerance, use Eq.~\eqref{eq:boundary-vp} rather than the generic circle-intersection map.
  \item \emph{Nearly collinear control points:} if the determinant in Eq.~\eqref{eq:circumcircle-determinant} is small relative to the coordinate scale, draw a straight segment.
  \item \emph{Circle tangency:} clamp a slightly negative value of $r_1^2-\ell^2$ to zero before taking the square root.
  \item \emph{Wrong intersection branch:} when both coordinate-circle intersections are real, enforce the image-disk condition and select the candidate nearest the center.
  \item \emph{Angular wrap-around:} normalize angles consistently and explicitly test which clockwise or counterclockwise interval contains the control point.
  \item \emph{Rendering order:} draw guide arcs behind the solid cube arcs so that the cube silhouette remains visually dominant.
\end{itemize}

\section{Conclusion}
Curvilinear perspective is most naturally understood by separating spherical geometry from planar rendering.  A spatial straight line always produces a great circle of viewing directions, and its two orientations determine a pair of vanishing points shared by all parallel lines.  Stereographic projection maps that great circle exactly to a circle or straight line.  The practical projection used in the project follows a different but closely related construction: horizontal and vertical viewing angles define intersecting circular coordinate lines, the transformed cube axes determine three vanishing-point pairs, and each finite cube edge is represented by a circumcircle constrained by its endpoints and its direction-family vanishing point.  This yields a controllable five-point-style curvilinear perspective with clear limiting cases, while also making explicit where circularity is exact and where it is an implementation-level approximation.
